\chapter{Introduction}\label{ch:intro}

\section{Motivation}
The battery pack is the most expensive and the most safety-critical part of an electric vehicle. A lithium-ion cell is safe only inside a narrow window of voltage, current and temperature; outside it the cell ages quickly, and at the extreme it can go into thermal runaway \citep{feng2018}. The battery management system (BMS) keeps the pack inside this window. It does three things at the same time: it measures what it can (cell voltages, pack current, a few temperatures), it estimates what it cannot measure (state of charge, core temperature, internal resistance), and it decides how much current the vehicle may draw or push back.

The estimation part matters because the dangerous quantities are hidden. The temperature that matters for degradation and for runaway is the temperature inside the cell, but a sensor can only be glued to the surface. The state of charge cannot be read off a voltmeter while current flows, because the terminal voltage also depends on the internal resistance and on slow polarisation effects. A BMS that only compares raw measurements against limits reacts late, and a BMS that trusts an estimate blindly reacts wrongly when the model does not match the cell.

The decision part matters because the obvious reaction to a problem, opening the main contactor, is itself dangerous in a moving car \citep{lelie2018}. A good BMS therefore reacts in steps: it first reduces the allowed current, then limits the car to a ``limp-home'' level, and only disconnects the pack when nothing else is safe.

Testing all this is the hard part. Many of the cases a BMS must handle (a broken temperature sensor, a failed cooling system, an external heat source, a lost measurement line, a software hang) cannot be produced on a real pack safely or repeatably. Industry solves this with hardware-in-the-loop (HIL) rigs that simulate the pack around a real BMS controller, but such rigs are expensive laboratory equipment \citep{barreras2016,haupt2013}.

\section{Aim and scope}
The aim of this thesis is to design a complete BMS controller and to verify it in closed loop against a simulated pack, on hardware that a student can afford. Concretely:
\begin{itemize}
  \item BMS software for a 19S12P pack of LG INR21700 M50-family cells, running in real time on an automotive-class microcontroller (STM32G431, Arm Cortex-M4F) and written in register-level C.
  \item A digital twin of the pack, parameterised with published measurements of the M50T cell \citep{kirkaldy2024}, running on a separate microcontroller and connected to the BMS only through a real CAN bus.
  \item A verification chain from model-in-the-loop to hardware-in-the-loop, including randomised Monte Carlo testing on the target microcontroller.
\end{itemize}
The HIL setup is \emph{communication-based}: the twin sends cell voltages, current and temperature as CAN messages instead of producing analog cell voltages. This tests the control strategy, which is the focus here; analog front-end hardware, insulation monitoring and contactor pre-charge are out of scope (\cref{ch:discussion}).

\section{Contributions}
\begin{enumerate}
  \item A working BMS controller on real hardware with a real CAN bus: SOC estimation, core-temperature estimation, thermal anomaly detection, sensor diagnostics, a graded decision layer and functional-safety mechanisms (watchdog, flash CRC, end-to-end protected messages). Only the input side would change for a real pack.
  \item A low-cost, reusable test platform: an Arduino Uno running the pack twin, an STM32 Nucleo board as the controller and two CAN transceivers. It runs faults that are impossible on real cells (60\,°C cells, a cut cable, a software hang) repeatably and from scripts.
  \item Evidence for the method: the HIL bench found three design flaws that the earlier model-, software- and processor-in-the-loop stages had not found.
  \item A model grounded in measured data: the open-circuit voltage, the resistances, their temperature dependence and the ageing relation are fitted to public measurements of the M50T cell, and the parameters that are still assumed are listed explicitly.
  \item A remote-monitoring path to a web dashboard, with measured link quality.
\end{enumerate}

\section{Structure of the thesis}
\Cref{ch:background} summarises the background and the related HIL literature. \Cref{ch:design} describes the hardware, the software structure and the CAN design. \Cref{ch:models} identifies and validates the cell model. \Cref{ch:estimation} presents the estimators and \cref{ch:decision} the decision layer and the stress index. \Cref{ch:verification} reports the verification results, \cref{ch:safety} the safety analysis and \cref{ch:monitoring} the remote monitoring. \Cref{ch:discussion} discusses the results and their limits, and \cref{ch:conclusion} concludes.
